\chapter{Frame Equivariance of the $s_1$ Tracker}
\label{app:frame_equivariance}
% Ported from: Draft_11.tex Sec. "Relation Between the Fields" (D is not
% objective; s1 is objective) and Sec. "Behavior Under a Rotating
% Observer" (the double-gyre rotating-frame experiment and its result).

This appendix collects the statement, and what supporting derivation
exists in the source paper, for why the $s_1$ tracker of
Chapter~\ref{ch:second_order} is objective (frame-indifferent) while the
$D$ tracker is not.

\section{The Determinant Field Is Not Objective}
\label{sec:appD:D_not_objective}

Both the determinant field $D=\det\mathbf{J}$ and the smaller strain
eigenvalue $s_1$ are read from the same fitted Jacobian $\hat{\mathbf{J}}$.
Where the vorticity $\omega = \partial v/\partial x - \partial u/\partial y$
vanishes, $\mathbf{J}$ equals its symmetric part $\mathbf{S}$ and
$D = s_1s_2$, which is $-s_1^2$ in incompressible flow. Along such a
curve the two fields share their trenches. Elsewhere the spin
$\mathbf{W} = \mathbf{J}-\mathbf{S}$, the antisymmetric part of
$\mathbf{J}$ that carries the vorticity, separates them, and it also
separates them across observers.

A frame rotating clockwise at rate $\Omega$, with rotation $\mathbf{Q}$,
adds a uniform spin that raises $\omega$ to $\omega+2\Omega$. The strain
tensor only co-rotates, so $s_1$ is objective \cite{52}, while the
determinant keeps the spin,
\begin{equation}
D' = D + \Omega\,\omega + \Omega^2.
\label{eq:appD:D_not_objective}
\end{equation}
Relatively rotating observers therefore disagree about where the
trenches of $D$ lie.
%% NOTE(new-math, checked 2026-09-29): Draft_11 states
%% eq:appD:D_not_objective without the algebra; the derivation below is new.

To see where Equation~\eqref{eq:appD:D_not_objective} comes from, split
the Jacobian into strain and spin,
\begin{equation}
\mathbf{S} = \mu\mathbf{I} + \begin{bmatrix} s_n & s_s \\ s_s & -s_n\end{bmatrix},
\qquad
\mathbf{W} = \frac{\omega}{2}\begin{bmatrix} 0 & -1 \\ 1 & 0\end{bmatrix},
\label{eq:appD:split}
\end{equation}
with $\mu = \tfrac12\operatorname{div}\mathbf{v}$ and
$r = \sqrt{s_n^2+s_s^2}$, so the strain eigenvalues are $s_{1,2} = \mu \mp r$.
Expanding the determinant,
\begin{equation}
D = \det(\mathbf{S}+\mathbf{W}) = \mu^2 - r^2 + \frac{\omega^2}{4}
  = s_1 s_2 + \frac{\omega^2}{4}.
\label{eq:appD:D_split}
\end{equation}
The rotating observer of Equation~\eqref{eq:appD:rotating_field} adds the
field $\Omega(-y',x')^{\mathsf T}$, whose gradient is purely antisymmetric,
so it leaves $\mathbf{S}$ unchanged up to the co-rotation
$\mathbf{S}' = \mathbf{Q}\mathbf{S}\mathbf{Q}^{\mathsf T}$, which preserves
$\mu$ and $r$, and replaces $\omega/2$ by $\omega/2 + \Omega$. Substituting
into Equation~\eqref{eq:appD:D_split},
\begin{equation}
D' = \mu^2 - r^2 + \Bigl(\frac{\omega}{2}+\Omega\Bigr)^2
   = D + \Omega\,\omega + \Omega^2 ,
\end{equation}
while $s_1' = s_1$ and $s_2' = s_2$. The $s_1$ tracker's gradient and
Hessian are those of an observer-invariant scalar field, and its tangent
reference $\mathbf{t}_{\mathrm{ref}}$ is carried in state rather than read
from a world axis, so its commanded velocity co-rotates with the frame.

\section{Behavior Under a Rotating Observer}
\label{sec:appD:rotating}

Both trackers were run from $(0.05,0.40)$, just inside the upper saddle
of the double gyre, on the same physical flow in two observer frames: the
inertial one and a frame rotating clockwise at $\Omega=0.23$, a heading
drift of about $8^\circ$/h at the ocean time unit of
Chapter~\ref{ch:second_order}, where the field observed in the rotating
frame is
\begin{equation}
\mathbf{v}'(\mathbf{p}',t) = \mathbf{Q}\mathbf{v}(\mathbf{Q}^\top\mathbf{p}') + \Omega\,(-y',x')^\top.
\label{eq:appD:rotating_field}
\end{equation}
The rate keeps the structure's apparent speed along the tracked path,
$\Omega\lVert\mathbf{p}\rVert \leq 0.11$, inside the command authority
$k\,c_{\max}=0.12$, the condition required by
Appendix~\ref{app:stability}. Rotating-frame trajectories are pulled back
to inertial coordinates and compared with the inertial run of the same
primitive.

The trajectories are shown in Figure~\ref{fig:ch7:oecs_objectivity}.

Between acquisition and the far saddle, the $s_1$ tracker stays within
0.016 of the separatrix in the rotating frame, against 0.010 in the
inertial frame, keeps the straddle throughout, and both runs capture at
the same bottom saddle. The $D$ tracker's rotating-frame run is displaced
by up to 0.084, against 0.010, and loses the straddle for 36 steps across
the crest, with all six robots on one side of the separatrix. Sweeping
the rate, the $D$ tracker first loses the straddle at
$\Omega \approx 0.21$, where $\Omega\lVert\mathbf{p}\rVert$ at the crest
is a sixth of the command authority.

The frame's solid-body vorticity shifts the $D'$ landscape
(Equation~\eqref{eq:appD:D_not_objective}) and adds a transport term to
the measured flow, so both terms of the $D$ tracker are steered by a
frame artifact. On the separatrix $\omega=0$ but
$\partial_x\omega = 2\pi^3A\sin\pi y_f$, so $\Omega\omega$ tilts $D'$
across the trench. Against the transverse curvature $2\pi^6A^2$, it moves
the trench by $\Omega\sin(\pi y_f)/(\pi^3A)$, largest at the origin and
zero at the saddles, so the run rejoins the separatrix near
$\mathbf{p}_1^*$.

\section{Why the $s_1$ Command Is Equivariant}
\label{sec:appD:s1_equivariant}

The $s_1$ tracker's tests are scalar inequalities in quantities a
rotation leaves invariant, so its command co-rotates with the frame, and
its residual error is the frame-transport term bounded in
Appendix~\ref{app:stability}. Its one non-objective input is the tangent
seed taken from the measured flow, which fixes the ride direction on the
first ride step (Equation~\eqref{eq:ch7:tangent_select}); the frames agree
while
\begin{equation}
\Omega\lVert\mathbf{p}_{\text{seed}}\rVert < \bigl\lvert\hat{\mathbf{v}}_0^\top\mathbf{t}_{\text{raw}}\bigr\rvert.
\label{eq:appD:seed_condition}
\end{equation}
Past that one seed step, the tangent is carried in state through
$\mathbf{t}_{\mathrm{ref}}$ (Equation~\eqref{eq:ch7:tangent_select}), which is
itself built from objective quantities, so the $s_1$ tracker is objective
except for this single, bounded, first-step dependence on the measured
flow.
%% NOTE: \ref{eq:ch7:tangent_select} matches the eq:ch7:-prefixed label
%% currently defined in ch07_second_order.tex for the $s_1$ tracker's
%% tangent-selection rule; re-check if that chapter's labels change.
